\section{Заключение}

В настоящей работе разработан программный комплекс для прогнозирования
реализованной волатильности рынка драгоценных металлов Московской биржи
с учётом настроений и внимания инвесторов.

\textbf{Основные результаты:}

\begin{enumerate}
    \item Собран и обработан панельный датасет дневных котировок золота (GLDRUB\_TOM)
          и серебра (SLVRUB\_TOM) за период 2018--2026 гг.\ (1\,931 и 1\,635 торговых дней
          соответственно). Реализованная волатильность вычислена оценщиком Garman--Klass.

    \item Сформирован многоисточниковый вектор sentiment/attention/macro (25 признаков):
          новостной поток RSS (6 источников), GDELT (718 дней с покрытием за 2018--2026),
          Telegram (12 каналов, 23\,733 сообщений, 2\,028 дней),
          Google Trends (3\,043 дня), курсы валют и ключевая ставка ЦБ РФ,
          объёмы торгов MOEX, макроэкономические данные FRED (Brent, DXY, VIX, CPI США, ставка ФРС).
          Все источники автоматически обновляются через реализованные парсеры.

    \item Эконометрический бенчмарк HAR-log показал устойчивые результаты:
          $R^2_{\mathrm{OOS}} = 0{,}177$ (золото) и $0{,}377$ (серебро).
          Коэффициенты $\hat\beta_m > \hat\beta_w > \hat\beta_d > 0$
          соответствуют теоретическим предсказаниям модели для золота.

    \item Расширенная спецификация HAR-log-S с лагированным взвешенным
          sentiment-индексом улучшает MAE на обоих инструментах
          ($0{,}674 \to 0{,}672$ для золота, $0{,}838 \to 0{,}834$ для серебра),
          однако улучшение значимо лишь на уровне 10\% по тесту
          Диболда--Мариано (раздел~\ref{sec:dm}).
          Ablation-анализ показал, что вклад различных sentiment-источников
          различается между инструментами: для золота сильнейший источник~---
          взвешенный \texttt{sentiment\_combined}, для серебра~---
          \texttt{attention\_google} (Google Trends).
          Это согласуется с теорией Да--Энгельберга--Гао~\cite{Da2011}
          об индексе поисковой активности как самостоятельной мере
          внимания розничных инвесторов.

    \item Модель XGBoost с расширенным sentiment-вектором даёт
          $R^2_{\mathrm{OOS}} = 0{,}070$ (золото) и $0{,}305$ (серебро),
          уступая HAR-log/HAR-log-S на обоих инструментах.

    \item kNN-модель на пространстве HAR+sentiment-признаков
          (Halousková \& Lyócsa,~\cite{Halouskova2025};
          $k^*=75$ для золота, $k^*=30$ для серебра)
          после устранения потенциальной утечки данных (лаг sentiment~1)
          и унификации формулы QLIKE \textit{не превзошла} HAR-log/HAR-log-S
          по $R^2_{\mathrm{OOS}}$ и QLIKE на обоих инструментах.
          Этот негативный результат, выявленный в ходе работы, важен сам по себе:
          он показывает, что прежние внутренние оценки преимущества kNN
          были артефактом несовместимых определений метрик и одновременности
          sentiment с целевой переменной.

    \item Реализован веб-интерфейс на Flask с тремя страницами (обзор,
          прогноз волатильности, sentiment-монитор) для интерактивной
          визуализации реализованной волатильности, прогноза модели
          и динамики sentiment/attention-индексов.
\end{enumerate}

\textbf{Методологический вклад работы:}
показано, что для строгого сравнения моделей прогнозирования
волатильности необходимо одновременное соблюдение двух условий~---
единой формулы QLIKE и согласованного временного лага всех экзогенных
признаков относительно целевой переменной. Нарушение любого из условий
может привести к иллюзорному преимуществу более сложных моделей.
Для проекта подготовлен единый модуль метрик (\texttt{src/models/metrics.py})
с канонической Patton QLIKE и тестом Диболда--Мариано,
который может быть переиспользован в дальнейших работах.
Основной практический итог проекта~--- не отдельная модель, а законченный
комплекс: сбор семи источников (MOEX ISS, GDELT, Google Trends, Telegram,
RSS, ЦБ~РФ, FRED), приведение их к единому календарю торгов, организация
оценки на расширяющемся окне без утечек и веб-визуализация результатов.
Именно объём инженерной части (раздел~\ref{sec:volume}) и составляет основную
сложность работы.

\textbf{Направления дальнейшей работы:}

\begin{itemize}
    \item Расширение NLP-покрытия RSS-новостей: текущая реализация
          даёт устойчиво нулевой sentiment для большинства торговых дней,
          что нейтрализует регрессор в ablation
          (таблица~\ref{tab:har_ablation}).
    \item Подключение дополнительных макро-сигналов
          (объёмы СВОП-сделок ЦБ РФ, индексы геополитической
          неопределённости~GPR).
    \item Применение методологии к платине и палладию по мере
          накопления выборки до уровня, достаточного для
          надёжной оценки expanding-window OOS
          (в текущем датасете 346 наблюдений на инструмент).
    \item Распространение HAR-log-S на более длинные горизонты
          прогнозирования (5, 22 дня) с агрегацией sentiment-вектора
          по соответствующим окнам.
\end{itemize}
